\section{Plan de decisiones técnicas de K2: cómo se condicionan entre sí los cuatro objetivos}

La sección anterior explicó que K2 no es una optimización aislada; esta sección organiza esos objetivos en un plan de decisiones y responde a la pregunta «¿qué optimiza a la vez una empresa de modelos?». El modelo base, las capacidades agénticas, el ecosistema de código abierto y la comercialización no son consignas paralelas, sino cuatro objetivos que se condicionan entre sí. Cuanto más fuerte es el modelo base, más respaldo tienen los productos de Agent y la API; cuanto más completa es la apertura del código, más rápido se difunde el ecosistema, pero más difícil resulta la captura comercial; cuanto más fuertes son las capacidades agénticas, más importan las herramientas y la retroalimentación, y más complejos se vuelven el entrenamiento y la evaluación.

\begin{knowledgebox}{Matriz de los cuatro objetivos de K2}
\begin{tabularx}{\textwidth}{p{0.18\textwidth}p{0.30\textwidth}p{0.38\textwidth}}
\toprule
Objetivo & Qué se optimiza & Qué arrastra consigo \\
\midrule
Base Model & Capacidad general, capacidad de programación, capacidad de razonamiento & Calidad de los datos, token efficiency, optimizador y capacidad de cómputo. \\
Agentic & Uso de herramientas en varios turnos, retroalimentación del entorno, generalización entre tareas & RL, cómputo en tiempo de prueba, protocolos de herramientas y evaluación. \\
Open Source & Validación externa, ecosistema de desarrolladores, confianza & Captura comercial, costo de soporte y reutilización por parte de competidores. \\
Business & API, productos, pagos de empresas y desarrolladores & Límites del modelo, límites de las aplicaciones de Agent y estructura de costos. \\
\bottomrule
\end{tabularx}
\end{knowledgebox}

\begin{importantbox}{Por qué estos cuatro objetivos no pueden analizarse por separado}
Si solo hay un modelo base fuerte sin capacidades agénticas, el modelo sigue siendo como un «cerebro en una cubeta»; si solo hay productos de Agent con un modelo base débil, la generalización queda limitada; si hay código abierto pero no una vía comercial, el ecosistema difícilmente se sostiene; si la comercialización llega demasiado pronto, puede limitar la investigación y la difusión del código abierto.
\end{importantbox}

\subsection{El ciclo entre eficiencia de entrenamiento y capacidad del producto}

En la matriz de los cuatro objetivos, lo que más fácilmente se subestima es la eficiencia de entrenamiento; esta subsección explica por qué se transmite al producto. La token efficiency parece una métrica de entrenamiento, pero influye en la capacidad del producto. Una mayor token efficiency significa obtener un modelo base más fuerte con los mismos datos y la misma capacidad de cómputo; un modelo base más fuerte permite que el Agent resuelva mejor tareas complejas; las tareas complejas generan más retroalimentación y más datos de uso; y esa retroalimentación puede incorporarse a la siguiente ronda de entrenamiento y de iteración del producto. Si este ciclo se pone en marcha, la empresa de modelos ya no solo entrena modelos, sino que construye un volante de inercia de capacidades.

\begin{knowledgebox}{Volante de inercia de capacidades}
\begin{tabularx}{\textwidth}{p{0.20\textwidth}p{0.34\textwidth}X}
\toprule
Etapa & Función & Riesgo \\
\midrule
Eficiencia de entrenamiento & Reducir el costo de obtener capacidades & Inestabilidad del optimizador o de la estrategia de datos. \\
Producto de Agent & Llevar el modelo a tareas reales & Tasa de fallas, seguridad de las herramientas y confianza del usuario. \\
Datos de retroalimentación & Registrar éxitos y fallas en tareas reales & Mucho ruido, privacidad y sesgo de selección. \\
Siguiente ronda de entrenamiento & Corregir el modelo y la estrategia de herramientas con la retroalimentación & Una evaluación imprecisa amplifica los errores. \\
\bottomrule
\end{tabularx}
\end{knowledgebox}

\subsection{Resumen de la sección}

Las decisiones técnicas de K2 no buscan el óptimo de un solo punto, sino la optimización conjunta de cuatro objetivos. Solo al entender este plan se entiende por qué la entrevista va y viene entre la token efficiency, la generalización agéntica, el código abierto, la API y el modelo de negocio.
